
\documentclass{article}
\usepackage{colortbl}
\usepackage{makecell}
\usepackage{multirow}
\usepackage{supertabular}

\begin{document}

\newcounter{utterance}

\centering \large Interaction Transcript for game `matchit\_ascii', experiment `different\_grid', episode 2 with Apertus{-}v1.5{-}8B.
\vspace{24pt}

{ \footnotesize  \setcounter{utterance}{1}
\setlength{\tabcolsep}{0pt}
\begin{supertabular}{c@{$\;$}|p{.15\linewidth}@{}p{.15\linewidth}p{.15\linewidth}p{.15\linewidth}p{.15\linewidth}p{.15\linewidth}}
    \# & $\;$A & \multicolumn{4}{c}{Game Master} & $\;\:$B\\
    \hline

    \theutterance \stepcounter{utterance}  
    & & \multicolumn{4}{p{0.6\linewidth}}{
        \cellcolor[rgb]{0.9,0.9,0.9}{
            \makecell[{{p{\linewidth}}}]{
                \texttt{\tiny{[P1$\langle$GM]}}
                \texttt{Učestvujete u zajedničkoj igri pogađanja.} \\
\texttt{Cilj je da utvrdimo da li su ova mreža i druga slika koju samo ja mogu da vidim iste.} \\
\texttt{Molim vas da prvo opišete svoju mrežu. Zatim ću ja dati svoj opis i možemo postavljati jedno drugom pitanja o mrežama kako bismo utvrdili da li su iste. Sada započnite svoj kratak opis mreže sa "OPIS:" a zatim navedite opis. Ne dodajte ništa drugo. Ovo je vaša mreža:} \\
\texttt{X X X X X} \\
\texttt{X ▢ ▢ ▢ ▢} \\
\texttt{X ▢ ▢ ▢ ▢} \\
\texttt{X ▢ ▢ ▢ ▢} \\
\texttt{X ▢ ▢ ▢ ▢} \\
            }
        }
    }
    & \\ \\

    \theutterance \stepcounter{utterance}  
    & \multicolumn{4}{p{0.6\linewidth}}{
        \cellcolor[rgb]{0.9,0.9,0.9}{
            \makecell[{{p{\linewidth}}}]{
                \texttt{\tiny{[P1$\rangle$GM]}}
                \texttt{OPIS: Mreža prikazuje pet redova horizontalnih elemenata. Prvi red sadrži pet X simbolova. Sva ostala četiri reda imaju isti uzorak: započinju jednim X simbolom, a zatim imaju četiri ▢ simbola redovno raspoređena.} \\
            }
        }
    }
    & & \\ \\

    \theutterance \stepcounter{utterance}  
    & & & \multicolumn{2}{p{0.3\linewidth}}{
        \cellcolor[rgb]{0.9,0.9,0.9}{
            \makecell[{{p{\linewidth}}}]{
                \texttt{\tiny{[GM$|$GM]}}
                \texttt{continue} \\
            }
        }
    }
    & & \\ \\

    \theutterance \stepcounter{utterance}  
    & & \multicolumn{4}{p{0.6\linewidth}}{
        \cellcolor[rgb]{0.9,0.9,0.9}{
            \makecell[{{p{\linewidth}}}]{
                \texttt{\tiny{[GM$\rangle$P2]}}
                \texttt{Učestvujete u zajedničkoj igri pogađanja.} \\
\texttt{Cilj je da utvrdimo da li su ova mreža i druga slika koju samo ja mogu da vidim iste.} \\
\texttt{Molim vas da prvo opišete svoju mrežu. Zatim ću ja dati svoj opis i možemo postavljati jedno drugom pitanja o mrežama kako bismo utvrdili da li su iste. Sada započnite svoj kratak opis mreže sa "OPIS:" a zatim navedite opis. Ne dodajte ništa drugo. Ovo je vaša mreža:} \\
\texttt{▢ ▢ X ▢ ▢} \\
\texttt{▢ ▢ X ▢ ▢} \\
\texttt{▢ X X X ▢} \\
\texttt{▢ X X X ▢} \\
\texttt{X X X X X} \\
            }
        }
    }
    & \\ \\

    \theutterance \stepcounter{utterance}  
    & & & \multicolumn{4}{p{0.6\linewidth}}{
        \cellcolor[rgb]{0.9,0.9,0.9}{
            \makecell[{{p{\linewidth}}}]{
                \texttt{\tiny{[GM$\langle$P2]}}
                \texttt{OPIS: Mreža se sastoji od pet redova i sedam stupaca. Na položajima (2,1), (2,3), (2,5) nalazi se znak X. Na položajima (3,2), (3,3), (3,4) nalaze se znakovi X. Na položajima (4,2), (4,3), (4,4) nalaze se znakovi X. Na položajima (5,1), (5,2), (5,3), (5,4), (5,5), (5,6), (5,7) nalaze se znakovi X.} \\
            }
        }
    }
     \\ \\

    \theutterance \stepcounter{utterance}  
    & & & \multicolumn{2}{p{0.3\linewidth}}{
        \cellcolor[rgb]{0.9,0.9,0.9}{
            \makecell[{{p{\linewidth}}}]{
                \texttt{\tiny{[GM$|$GM]}}
                \texttt{continue} \\
            }
        }
    }
    & & \\ \\

    \theutterance \stepcounter{utterance}  
    & & \multicolumn{4}{p{0.6\linewidth}}{
        \cellcolor[rgb]{0.9,0.9,0.9}{
            \makecell[{{p{\linewidth}}}]{
                \texttt{\tiny{[GM$\rangle$P2]}}
                \texttt{Ovo je mojOPIS: Mreža prikazuje pet redova horizontalnih elemenata. Prvi red sadrži pet X simbolova. Sva ostala četiri reda imaju isti uzorak: započinju jednim X simbolom, a zatim imaju četiri ▢ simbola redovno raspoređena.} \\
\texttt{Sada postavi pitanje kako bi saznao nove aspekte moje mreže koji se mogu razlikovati od tvoje mreže. Počni sa "PITANJE:" i ne dodaj ništa drugo.} \\
            }
        }
    }
    & \\ \\

    \theutterance \stepcounter{utterance}  
    & & & \multicolumn{4}{p{0.6\linewidth}}{
        \cellcolor[rgb]{0.9,0.9,0.9}{
            \makecell[{{p{\linewidth}}}]{
                \texttt{\tiny{[GM$\langle$P2]}}
                \texttt{PITANJE: Da li je svaki red u vašoj mreži istog dužine, tj. istog broja elemenata?} \\
            }
        }
    }
     \\ \\

    \theutterance \stepcounter{utterance}  
    & & & \multicolumn{2}{p{0.3\linewidth}}{
        \cellcolor[rgb]{0.9,0.9,0.9}{
            \makecell[{{p{\linewidth}}}]{
                \texttt{\tiny{[GM$|$GM]}}
                \texttt{continue} \\
            }
        }
    }
    & & \\ \\

    \theutterance \stepcounter{utterance}  
    & & \multicolumn{4}{p{0.6\linewidth}}{
        \cellcolor[rgb]{0.9,0.9,0.9}{
            \makecell[{{p{\linewidth}}}]{
                \texttt{\tiny{[P1$\langle$GM]}}
                \texttt{Ovo je mojOPIS: Mreža se sastoji od pet redova i sedam stupaca. Na položajima (2,1), (2,3), (2,5) nalazi se znak X. Na položajima (3,2), (3,3), (3,4) nalaze se znakovi X. Na položajima (4,2), (4,3), (4,4) nalaze se znakovi X. Na položajima (5,1), (5,2), (5,3), (5,4), (5,5), (5,6), (5,7) nalaze se znakovi X.} \\
\texttt{PITANJE: Da li je svaki red u vašoj mreži istog dužine, tj. istog broja elemenata?Započnite svoj odgovor sa "ODGOVOR:" i ne dodajte ništa drugo.} \\
            }
        }
    }
    & \\ \\

    \theutterance \stepcounter{utterance}  
    & \multicolumn{4}{p{0.6\linewidth}}{
        \cellcolor[rgb]{0.9,0.9,0.9}{
            \makecell[{{p{\linewidth}}}]{
                \texttt{\tiny{[P1$\rangle$GM]}}
                \texttt{ODGOVOR: Da, svaki red u mreži ima isti broj elemenata, odnosno istu dužinu. Postoje ukupno pet redova, a svaki red sadrži sedam elemenata.} \\
            }
        }
    }
    & & \\ \\

    \theutterance \stepcounter{utterance}  
    & & & \multicolumn{2}{p{0.3\linewidth}}{
        \cellcolor[rgb]{0.9,0.9,0.9}{
            \makecell[{{p{\linewidth}}}]{
                \texttt{\tiny{[GM$|$GM]}}
                \texttt{continue} \\
            }
        }
    }
    & & \\ \\

    \theutterance \stepcounter{utterance}  
    & & \multicolumn{4}{p{0.6\linewidth}}{
        \cellcolor[rgb]{0.9,0.9,0.9}{
            \makecell[{{p{\linewidth}}}]{
                \texttt{\tiny{[P1$\langle$GM]}}
                \texttt{Sada postavi pitanje kako bi saznao nove aspekte moje mreže koji se mogu razlikovati od tvoje mreže. Počni sa "PITANJE:" i ne dodaj ništa drugo.} \\
            }
        }
    }
    & \\ \\

    \theutterance \stepcounter{utterance}  
    & \multicolumn{4}{p{0.6\linewidth}}{
        \cellcolor[rgb]{0.9,0.9,0.9}{
            \makecell[{{p{\linewidth}}}]{
                \texttt{\tiny{[P1$\rangle$GM]}}
                \texttt{PITANJE: Da li su svaki ▢ element u vašoj mreži isti, ili postoji neka razlika između njih?} \\
            }
        }
    }
    & & \\ \\

    \theutterance \stepcounter{utterance}  
    & & & \multicolumn{2}{p{0.3\linewidth}}{
        \cellcolor[rgb]{0.9,0.9,0.9}{
            \makecell[{{p{\linewidth}}}]{
                \texttt{\tiny{[GM$|$GM]}}
                \texttt{continue} \\
            }
        }
    }
    & & \\ \\

    \theutterance \stepcounter{utterance}  
    & & \multicolumn{4}{p{0.6\linewidth}}{
        \cellcolor[rgb]{0.9,0.9,0.9}{
            \makecell[{{p{\linewidth}}}]{
                \texttt{\tiny{[GM$\rangle$P2]}}
                \texttt{ODGOVOR: Da, svaki red u mreži ima isti broj elemenata, odnosno istu dužinu. Postoje ukupno pet redova, a svaki red sadrži sedam elemenata.} \\
\texttt{PITANJE: Da li su svaki ▢ element u vašoj mreži isti, ili postoji neka razlika između njih?Započnite svoj odgovor sa "ODGOVOR:" i ne dodajte ništa drugo.} \\
            }
        }
    }
    & \\ \\

    \theutterance \stepcounter{utterance}  
    & & & \multicolumn{4}{p{0.6\linewidth}}{
        \cellcolor[rgb]{0.9,0.9,0.9}{
            \makecell[{{p{\linewidth}}}]{
                \texttt{\tiny{[GM$\langle$P2]}}
                \texttt{ODGOVOR: Sve ▢ elementi su identični u mreži; nema razlike između njih.} \\
            }
        }
    }
     \\ \\

    \theutterance \stepcounter{utterance}  
    & & & \multicolumn{2}{p{0.3\linewidth}}{
        \cellcolor[rgb]{0.9,0.9,0.9}{
            \makecell[{{p{\linewidth}}}]{
                \texttt{\tiny{[GM$|$GM]}}
                \texttt{continue} \\
            }
        }
    }
    & & \\ \\

    \theutterance \stepcounter{utterance}  
    & & \multicolumn{4}{p{0.6\linewidth}}{
        \cellcolor[rgb]{0.9,0.9,0.9}{
            \makecell[{{p{\linewidth}}}]{
                \texttt{\tiny{[GM$\rangle$P2]}}
                \texttt{Sada postavi pitanje kako bi saznao nove aspekte moje mreže koji se mogu razlikovati od tvoje mreže. Počni sa "PITANJE:" i ne dodaj ništa drugo.} \\
            }
        }
    }
    & \\ \\

    \theutterance \stepcounter{utterance}  
    & & & \multicolumn{4}{p{0.6\linewidth}}{
        \cellcolor[rgb]{0.9,0.9,0.9}{
            \makecell[{{p{\linewidth}}}]{
                \texttt{\tiny{[GM$\langle$P2]}}
                \texttt{PITANJE: Da li se X simboli pojavljuju na nekim identičnim položajima u vašoj mreži?} \\
            }
        }
    }
     \\ \\

    \theutterance \stepcounter{utterance}  
    & & & \multicolumn{2}{p{0.3\linewidth}}{
        \cellcolor[rgb]{0.9,0.9,0.9}{
            \makecell[{{p{\linewidth}}}]{
                \texttt{\tiny{[GM$|$GM]}}
                \texttt{continue} \\
            }
        }
    }
    & & \\ \\

    \theutterance \stepcounter{utterance}  
    & & \multicolumn{4}{p{0.6\linewidth}}{
        \cellcolor[rgb]{0.9,0.9,0.9}{
            \makecell[{{p{\linewidth}}}]{
                \texttt{\tiny{[P1$\langle$GM]}}
                \texttt{ODGOVOR: Sve ▢ elementi su identični u mreži; nema razlike između njih.} \\
\texttt{PITANJE: Da li se X simboli pojavljuju na nekim identičnim položajima u vašoj mreži?Započnite svoj odgovor sa "ODGOVOR:" i ne dodajte ništa drugo.} \\
            }
        }
    }
    & \\ \\

    \theutterance \stepcounter{utterance}  
    & \multicolumn{4}{p{0.6\linewidth}}{
        \cellcolor[rgb]{0.9,0.9,0.9}{
            \makecell[{{p{\linewidth}}}]{
                \texttt{\tiny{[P1$\rangle$GM]}}
                \texttt{ODGOVOR: Da, X simboli se pojavljuju na identičnim položajima u odnosu na položaje ▢ simbola u mreži.} \\
            }
        }
    }
    & & \\ \\

    \theutterance \stepcounter{utterance}  
    & & & \multicolumn{2}{p{0.3\linewidth}}{
        \cellcolor[rgb]{0.9,0.9,0.9}{
            \makecell[{{p{\linewidth}}}]{
                \texttt{\tiny{[GM$|$GM]}}
                \texttt{continue} \\
            }
        }
    }
    & & \\ \\

    \theutterance \stepcounter{utterance}  
    & & \multicolumn{4}{p{0.6\linewidth}}{
        \cellcolor[rgb]{0.9,0.9,0.9}{
            \makecell[{{p{\linewidth}}}]{
                \texttt{\tiny{[P1$\langle$GM]}}
                \texttt{Sada postavi pitanje kako bi saznao nove aspekte moje mreže koji se mogu razlikovati od tvoje mreže. Počni sa "PITANJE:" i ne dodaj ništa drugo.} \\
            }
        }
    }
    & \\ \\

    \theutterance \stepcounter{utterance}  
    & \multicolumn{4}{p{0.6\linewidth}}{
        \cellcolor[rgb]{0.9,0.9,0.9}{
            \makecell[{{p{\linewidth}}}]{
                \texttt{\tiny{[P1$\rangle$GM]}}
                \texttt{PITANJE: Da li postoje bilo kakvi razlici u uzorku položaja X simbola između vaše mreže i uzorka koji sam opisao u svom prvom opisu?} \\
            }
        }
    }
    & & \\ \\

    \theutterance \stepcounter{utterance}  
    & & & \multicolumn{2}{p{0.3\linewidth}}{
        \cellcolor[rgb]{0.9,0.9,0.9}{
            \makecell[{{p{\linewidth}}}]{
                \texttt{\tiny{[GM$|$GM]}}
                \texttt{continue} \\
            }
        }
    }
    & & \\ \\

    \theutterance \stepcounter{utterance}  
    & & \multicolumn{4}{p{0.6\linewidth}}{
        \cellcolor[rgb]{0.9,0.9,0.9}{
            \makecell[{{p{\linewidth}}}]{
                \texttt{\tiny{[GM$\rangle$P2]}}
                \texttt{ODGOVOR: Da, X simboli se pojavljuju na identičnim položajima u odnosu na položaje ▢ simbola u mreži.} \\
\texttt{PITANJE: Da li postoje bilo kakvi razlici u uzorku položaja X simbola između vaše mreže i uzorka koji sam opisao u svom prvom opisu?Započnite svoj odgovor sa "ODGOVOR:" i ne dodajte ništa drugo.} \\
            }
        }
    }
    & \\ \\

    \theutterance \stepcounter{utterance}  
    & & & \multicolumn{4}{p{0.6\linewidth}}{
        \cellcolor[rgb]{0.9,0.9,0.9}{
            \makecell[{{p{\linewidth}}}]{
                \texttt{\tiny{[GM$\langle$P2]}}
                \texttt{ODGOVOR: U odnosu na vaš opis, X simboli se redovno pojavljuju na svakom trećem položaju redova, a počinju u drugom položaju svakog reda izuzev prvog reda u kojem se X simboli pojavljuju na svim položajima.} \\
            }
        }
    }
     \\ \\

    \theutterance \stepcounter{utterance}  
    & & & \multicolumn{2}{p{0.3\linewidth}}{
        \cellcolor[rgb]{0.9,0.9,0.9}{
            \makecell[{{p{\linewidth}}}]{
                \texttt{\tiny{[GM$|$GM]}}
                \texttt{continue} \\
            }
        }
    }
    & & \\ \\

    \theutterance \stepcounter{utterance}  
    & & \multicolumn{4}{p{0.6\linewidth}}{
        \cellcolor[rgb]{0.9,0.9,0.9}{
            \makecell[{{p{\linewidth}}}]{
                \texttt{\tiny{[GM$\rangle$P2]}}
                \texttt{Sada postavi pitanje kako bi saznao nove aspekte moje mreže koji se mogu razlikovati od tvoje mreže. Počni sa "PITANJE:" i ne dodaj ništa drugo.} \\
            }
        }
    }
    & \\ \\

    \theutterance \stepcounter{utterance}  
    & & & \multicolumn{4}{p{0.6\linewidth}}{
        \cellcolor[rgb]{0.9,0.9,0.9}{
            \makecell[{{p{\linewidth}}}]{
                \texttt{\tiny{[GM$\langle$P2]}}
                \texttt{PITANJE: Da li postoji red u kojem se pojavljuju dva X simbola zaredom?} \\
            }
        }
    }
     \\ \\

    \theutterance \stepcounter{utterance}  
    & & & \multicolumn{2}{p{0.3\linewidth}}{
        \cellcolor[rgb]{0.9,0.9,0.9}{
            \makecell[{{p{\linewidth}}}]{
                \texttt{\tiny{[GM$|$GM]}}
                \texttt{continue} \\
            }
        }
    }
    & & \\ \\

    \theutterance \stepcounter{utterance}  
    & & \multicolumn{4}{p{0.6\linewidth}}{
        \cellcolor[rgb]{0.9,0.9,0.9}{
            \makecell[{{p{\linewidth}}}]{
                \texttt{\tiny{[P1$\langle$GM]}}
                \texttt{ODGOVOR: U odnosu na vaš opis, X simboli se redovno pojavljuju na svakom trećem položaju redova, a počinju u drugom položaju svakog reda izuzev prvog reda u kojem se X simboli pojavljuju na svim položajima.} \\
\texttt{PITANJE: Da li postoji red u kojem se pojavljuju dva X simbola zaredom?Započnite svoj odgovor sa "ODGOVOR:" i ne dodajte ništa drugo.} \\
            }
        }
    }
    & \\ \\

    \theutterance \stepcounter{utterance}  
    & \multicolumn{4}{p{0.6\linewidth}}{
        \cellcolor[rgb]{0.9,0.9,0.9}{
            \makecell[{{p{\linewidth}}}]{
                \texttt{\tiny{[P1$\rangle$GM]}}
                \texttt{ODGOVOR: Ne, u mreži ne postoji red u kojem se pojavljuju dva X simbola zaredom.} \\
            }
        }
    }
    & & \\ \\

    \theutterance \stepcounter{utterance}  
    & & & \multicolumn{2}{p{0.3\linewidth}}{
        \cellcolor[rgb]{0.9,0.9,0.9}{
            \makecell[{{p{\linewidth}}}]{
                \texttt{\tiny{[GM$|$GM]}}
                \texttt{continue} \\
            }
        }
    }
    & & \\ \\

    \theutterance \stepcounter{utterance}  
    & & \multicolumn{4}{p{0.6\linewidth}}{
        \cellcolor[rgb]{0.9,0.9,0.9}{
            \makecell[{{p{\linewidth}}}]{
                \texttt{\tiny{[P1$\langle$GM]}}
                \texttt{Sada postavi pitanje kako bi saznao nove aspekte moje mreže koji se mogu razlikovati od tvoje mreže. Počni sa "PITANJE:" i ne dodaj ništa drugo.} \\
            }
        }
    }
    & \\ \\

    \theutterance \stepcounter{utterance}  
    & \multicolumn{4}{p{0.6\linewidth}}{
        \cellcolor[rgb]{0.9,0.9,0.9}{
            \makecell[{{p{\linewidth}}}]{
                \texttt{\tiny{[P1$\rangle$GM]}}
                \texttt{PITANJE: Da li postoji bilo kakav red koji sadrži samo X simbole, bez ikakvog ▢ elementa?} \\
            }
        }
    }
    & & \\ \\

    \theutterance \stepcounter{utterance}  
    & & & \multicolumn{2}{p{0.3\linewidth}}{
        \cellcolor[rgb]{0.9,0.9,0.9}{
            \makecell[{{p{\linewidth}}}]{
                \texttt{\tiny{[GM$|$GM]}}
                \texttt{continue} \\
            }
        }
    }
    & & \\ \\

    \theutterance \stepcounter{utterance}  
    & & \multicolumn{4}{p{0.6\linewidth}}{
        \cellcolor[rgb]{0.9,0.9,0.9}{
            \makecell[{{p{\linewidth}}}]{
                \texttt{\tiny{[GM$\rangle$P2]}}
                \texttt{ODGOVOR: Ne, u mreži ne postoji red u kojem se pojavljuju dva X simbola zaredom.} \\
\texttt{PITANJE: Da li postoji bilo kakav red koji sadrži samo X simbole, bez ikakvog ▢ elementa?Započnite svoj odgovor sa "ODGOVOR:" i ne dodajte ništa drugo.} \\
            }
        }
    }
    & \\ \\

    \theutterance \stepcounter{utterance}  
    & & & \multicolumn{4}{p{0.6\linewidth}}{
        \cellcolor[rgb]{0.9,0.9,0.9}{
            \makecell[{{p{\linewidth}}}]{
                \texttt{\tiny{[GM$\langle$P2]}}
                \texttt{ODGOVOR: Da, prvi red sadrži samo X simbole, bez ikakvog ▢ elementa.} \\
            }
        }
    }
     \\ \\

    \theutterance \stepcounter{utterance}  
    & & & \multicolumn{2}{p{0.3\linewidth}}{
        \cellcolor[rgb]{0.9,0.9,0.9}{
            \makecell[{{p{\linewidth}}}]{
                \texttt{\tiny{[GM$|$GM]}}
                \texttt{continue} \\
            }
        }
    }
    & & \\ \\

    \theutterance \stepcounter{utterance}  
    & & \multicolumn{4}{p{0.6\linewidth}}{
        \cellcolor[rgb]{0.9,0.9,0.9}{
            \makecell[{{p{\linewidth}}}]{
                \texttt{\tiny{[GM$\rangle$P2]}}
                \texttt{Sada donesi odluku. Šta misliš: da li su tvoja mreža i druga opisana mreža iste?  Napiši "ODLUKA: ista mreža" ako misliš da su ista mreža ili "ODLUKA: različite mreže" ako misliš da su različite mreže. Ne dodaj ništa drugo.} \\
            }
        }
    }
    & \\ \\

    \theutterance \stepcounter{utterance}  
    & & & \multicolumn{4}{p{0.6\linewidth}}{
        \cellcolor[rgb]{0.9,0.9,0.9}{
            \makecell[{{p{\linewidth}}}]{
                \texttt{\tiny{[GM$\langle$P2]}}
                \texttt{ODLUKA: različite mreže} \\
            }
        }
    }
     \\ \\

    \theutterance \stepcounter{utterance}  
    & & & \multicolumn{2}{p{0.3\linewidth}}{
        \cellcolor[rgb]{0.9,0.9,0.9}{
            \makecell[{{p{\linewidth}}}]{
                \texttt{\tiny{[GM$|$GM]}}
                \texttt{continue} \\
            }
        }
    }
    & & \\ \\

    \theutterance \stepcounter{utterance}  
    & & & \multicolumn{2}{p{0.3\linewidth}}{
        \cellcolor[rgb]{0.9,0.9,0.9}{
            \makecell[{{p{\linewidth}}}]{
                \texttt{\tiny{[GM$|$GM]}}
                \texttt{success} \\
            }
        }
    }
    & & \\ \\

    \theutterance \stepcounter{utterance}  
    & & \multicolumn{4}{p{0.6\linewidth}}{
        \cellcolor[rgb]{0.9,0.9,0.9}{
            \makecell[{{p{\linewidth}}}]{
                \texttt{\tiny{[P1$\langle$GM]}}
                \texttt{ODGOVOR: Da, prvi red sadrži samo X simbole, bez ikakvog ▢ elementa.} \\
\texttt{Sada donesi odluku. Šta misliš: da li su tvoja mreža i druga opisana mreža iste?  Napiši "ODLUKA: ista mreža" ako misliš da su ista mreža ili "ODLUKA: različite mreže" ako misliš da su različite mreže. Ne dodaj ništa drugo.} \\
            }
        }
    }
    & \\ \\

    \theutterance \stepcounter{utterance}  
    & \multicolumn{4}{p{0.6\linewidth}}{
        \cellcolor[rgb]{0.9,0.9,0.9}{
            \makecell[{{p{\linewidth}}}]{
                \texttt{\tiny{[P1$\rangle$GM]}}
                \texttt{ODLUKA: različite mreže} \\
            }
        }
    }
    & & \\ \\

    \theutterance \stepcounter{utterance}  
    & & & \multicolumn{2}{p{0.3\linewidth}}{
        \cellcolor[rgb]{0.9,0.9,0.9}{
            \makecell[{{p{\linewidth}}}]{
                \texttt{\tiny{[GM$|$GM]}}
                \texttt{continue} \\
            }
        }
    }
    & & \\ \\

    \theutterance \stepcounter{utterance}  
    & & & \multicolumn{2}{p{0.3\linewidth}}{
        \cellcolor[rgb]{0.9,0.9,0.9}{
            \makecell[{{p{\linewidth}}}]{
                \texttt{\tiny{[GM$|$GM]}}
                \texttt{success} \\
            }
        }
    }
    & & \\ \\

\end{supertabular}
}

\end{document}
